\section*{Chapitre \ref{ch:qm:volatility-models} --- Modèles de volatilité}

\begin{solution}{exo:qm:volatility-models:1}
$\bar h = 0.02/(1 - 0.98) = 1$~: volatilité quotidienne 1~\%, annuelle $\sqrt{252} = 15.9\%$~; \omterm{def:qm:stochastic-differential-equations:ou}{demi-vie} $\ln\frac12/\ln0.98 = 34.3$ jours.
\end{solution}

\begin{solution}{exo:qm:volatility-models:2}
$\ln\frac12/\ln0.94 = 11.2$ jours. Les 20 derniers jours portent $1 - 0.94^{20} = 71\%$ du poids.
\end{solution}

\begin{solution}{exo:qm:volatility-models:3}
$m = 24 \times 12 = 288$ et $\sqrt{2/288} = 8.3\%$.
\end{solution}

\begin{solution}{exo:qm:volatility-models:4}
$h_{t+1} = 4\bar h$ et $\E_t[h_{t+k}] = \bar h + 0.98^{k-1} \times 3\bar h$~; en moyennant sur $k = 1, \dots, 10$ on obtient $3.74\bar h$, soit une volatilité de $\sqrt{3.74} = 1.93\%$
par jour. La prévision revient à peine vers la moyenne en dix jours, alors que la \omterm{def:qm:stochastic-differential-equations:ou}{demi-vie} est de 34 jours.
\end{solution}

\begin{solution}{exo:qm:volatility-models:5}
$\E[\hat\sigma^2/h + \ln h] = \sigma^2/h + \ln h$, dont la dérivée $-\sigma^2/h^2 + 1/h$ s'annule en $h = \sigma^2$, qui est un minimum. Avec $\sigma^2 = 1$, la perte vaut 1 à la
valeur vraie, $2 + \ln0.5 = 1.31$ pour $0.5\sigma^2$ et $0.67 + \ln1.5 = 1.07$ pour $1.5\sigma^2$~: une sous-prévision est davantage pénalisée.
\end{solution}

\begin{solution}{exo:qm:volatility-models:6}
Avec le rendement de 3.5~\%, la variance (à moyenne à peu près nulle) vaut $(249 \times 0.25 + 12.25)/250 = 0.298$, soit une volatilité de 0.546~\%~; sans lui, 0.5~\%. La
volatilité de la fenêtre baisse de 8.4~\% le jour où le rendement sort, quoi qu'il arrive ce jour-là.
\end{solution}

\begin{solution}{exo:qm:volatility-models:7}
Avec la vraisemblance gaussienne~: $\alpha = 0.029$, avec une \omterm{def:qm:estimation:estimator}{erreur type} hessienne de 0.0026 et une \omterm{def:qm:estimation:estimator}{erreur type} sandwich de 0.0033~; $\beta = 0.968$, avec 0.0028 et 0.0035. L'\omterm{def:qm:estimation:estimator}{erreur type} sandwich est environ 25~\%
plus grande parce que les résidus standardisés ont des queues épaisses, si bien que la vraisemblance gaussienne est mal spécifiée et que l'égalité de l'information échoue~; il faut donc rapporter l'\omterm{def:qm:estimation:estimator}{erreur type} sandwich.
Avec la vraisemblance de Student-$t$, les deux concordent (0.0033 et 0.0032 pour $\alpha$).
\end{solution}

\begin{solution}{exo:qm:volatility-models:8}
La valeur absolue du rendement est un proxy biaisé de la volatilité ($\E|r| = \sqrt{2/\pi}\,\sigma$ pour des rendements gaussiens, moins pour des queues épaisses), et l'erreur absolue moyenne n'appartient pas
à la famille robuste de Patton~: sous un proxy bruité, elle peut classer une prévision biaisée au-dessus de la variance vraie, en favorisant les prévisions trop basses. Comparer les prévisions de variance
avec l'erreur quadratique ou la \omterm{def:qm:volatility-models:eval}{perte QLIKE}, contre les rendements au carré ou la \omterm{def:qm:volatility-models:rv}{variance réalisée}, et tester la différence.
\end{solution}

\begin{solution}{pb:qm:volatility-models:1}
\textbf{1.} $\omega = 0.0010$, $\alpha = 0.029$ (0.003), $\beta = 0.968$ (0.004).
\textbf{2.} $\omega = 0.0007$, $\alpha = 0.030$ (0.003), $\beta = 0.968$ (0.003), $\nu = 7.2$ (0.6)~; log-vraisemblance supérieure de 141.
\textbf{3.} Persistance 0.9984~; \omterm{def:qm:stochastic-differential-equations:ou}{demi-vie} de 445 jours de bourse (255 avec la vraisemblance gaussienne).
\textbf{4.} $\sqrt{\omega/(1 - \alpha - \beta)} = 0.67\%$, contre un écart-type empirique de 0.58~\%~: avec une persistance aussi proche de un, le niveau implicite est mal
déterminé.
\textbf{5.} Non~: l'asymétrie GJR vaut 0.006, avec une \omterm{def:qm:estimation:estimator}{erreur type} de 0.005.
\textbf{6.} La hausse de 3.5~\% du 11 novembre 2022, qui quitte la fenêtre à la clôture du 2 novembre 2023.
\textbf{7.} De 0.538~\% à 0.496~\%, soit une baisse de 7.8~\%, alors que l'euro a gagné 1.2~\% ce jour-là.
\textbf{8.} Le GARCH est passé de 0.424~\% à 0.465~\%, l'EWMA de 0.410~\% à 0.490~\%.
\textbf{9.} Sous GARCH, son effet se réduit de moitié en 445 jours~; sous EWMA, en 11 jours~; dans la fenêtre, il reste entier pendant 250 jours et disparaît d'un coup.
\textbf{10.} Ce risque a diminué, un jour de mouvement d'une ampleur exceptionnelle.
\textbf{11.} QLIKE $-0.521$ (GARCH-$t$ estimé sur 1999--2014), $-0.495$ (EWMA), $-0.435$ (fenêtre), sur 3\,002 jours.
\textbf{12.} Oui~: \omterm{def:qm:volatility-models:eval}{test de Diebold--Mariano} $-2.76$ (GARCH contre EWMA, $p = 0.006$), $-4.10$ (GARCH contre la fenêtre), $-2.48$ (EWMA contre la fenêtre).
\textbf{13.} 0.77, 0.60 et 0.56~: tous inférieurs à un, donc chaque prévision réagit trop fortement aux données récentes~; le GARCH le moins.
\textbf{14.} Le proxy, un rendement quotidien au carré, est la variance multipliée par une variable de type $\chi^2_1$, si bien que l'essentiel de sa variation est du bruit qu'aucune prévision ne peut expliquer~; un $R^2$ de 3~\% est
proche du plafond.
\textbf{15.} Un proxy beaucoup moins bruité~: des classements plus nets, un $R^2$ plus élevé, et une puissance accrue dans les \omterm{def:qm:volatility-models:eval}{tests de Diebold--Mariano}.
\textbf{16.} La prévision GARCH-$t$, ou l'EWMA si la simplicité compte davantage que le gain de 5~\% en QLIKE~; jamais la fenêtre à pondérations égales.
\textbf{17.} Comme description de la lenteur avec laquelle les grands chocs s'estompent, oui~; comme valeur précise, non~: une persistance proche de un est la façon qu'a le modèle à mémoire courte d'approcher
la \omterm{def:qm:linear-time-series:longmem}{mémoire longue}, et l'intervalle de confiance de la \omterm{def:qm:stochastic-differential-equations:ou}{demi-vie} est large.
\textbf{18.} Que la baisse du risque publié était un artefact du calendrier de la fenêtre, et que toute prévision qui pondère selon le temps a monté.
\textbf{19.} Résultat nommé~: \emph{le fantôme du rapport de risque}~: un choc de variance EUR/USD a une \omterm{def:qm:stochastic-differential-equations:ou}{demi-vie} de 445 jours de bourse sous un GARCH(1,1) à innovations de Student-$t$~;
la volatilité de la fenêtre de 250 jours a baissé de 7.8~\% le 2 novembre 2023, le jour où le rendement de 3.5~\% du 11 novembre 2022 l'a quittée.
\textbf{20.} Elle garde chaque rendement à plein poids pendant une durée fixe, puis le retire, si bien que la prévision bouge au rythme des anniversaires plutôt qu'au gré des nouvelles.
\end{solution}

\begin{solution}{iq:qm:volatility-models:1}
Parce qu'un grand rendement quitte la fenêtre ce jour-là~: la baisse est l'anniversaire d'une vieille nouvelle, et non une information sur aujourd'hui, et elle peut l'emporter sur un grand rendement nouveau.
\iqlookfor{L'effet de sortie et sa correction (pondérations décroissantes).}
\end{solution}

\begin{solution}{iq:qm:volatility-models:2}
$h_t = \omega + \alpha r_{t-1}^2 + \beta h_{t-1}$~; variance inconditionnelle $\omega/(1 - \alpha - \beta)$ lorsque $\alpha + \beta < 1$~; $\alpha + \beta$ est la persistance, le taux de décroissance quotidien
d'un choc de variance, de \omterm{def:qm:stochastic-differential-equations:ou}{demi-vie} $\ln\frac12/\ln(\alpha + \beta)$. $\alpha$ est la réaction aux nouvelles, $\beta$ la mémoire.
\iqlookfor{La formule, la stationnarité et la persistance comme taux de décroissance.}
\end{solution}

\begin{solution}{iq:qm:volatility-models:3}
L'EWMA est un GARCH(1,1) avec $\omega = 0$ et $\alpha + \beta = 1$~: il est intégré. Il prévoit la variance d'aujourd'hui à tous les horizons, alors que le GARCH revient à son niveau de long terme~; pour
les longs horizons après un choc, l'EWMA surestime (ou, après une période calme, sous-estime) la variance.
\iqlookfor{L'IGARCH et la structure par terme plate des prévisions.}
\end{solution}

\begin{solution}{iq:qm:volatility-models:4}
Comparer à un proxy sans biais (rendements au carré ou \omterm{def:qm:volatility-models:rv}{variance réalisée}), avec une perte qui classe correctement sous le bruit du proxy, erreur quadratique ou QLIKE, tester la
différence des pertes avec le \omterm{def:qm:volatility-models:eval}{test de Diebold--Mariano}, et vérifier la calibration par une \omterm{def:qm:volatility-models:eval}{régression de Mincer--Zarnowitz}.
\iqlookfor{Proxy, perte robuste et test, et non un nombre unique.}
\end{solution}

\begin{solution}{iq:qm:volatility-models:5}
La somme des rendements intrajournaliers au carré, qui converge vers la \omterm{def:qm:volatility-models:rv}{variance intégrée} à mesure que l'échantillonnage devient plus fin. À très haute fréquence, le rebond bid-ask et la
discrétisation des prix ajoutent du bruit dont la contribution croît avec le nombre de rendements, si bien que la RV explose~; échantillonner toutes les quelques minutes, ou utiliser des \omterm{def:qm:estimation:estimator}{estimateurs} robustes au bruit
(chapitre~21).
\iqlookfor{La convergence, et le compromis avec le bruit.}
\end{solution}

\begin{solution}{iq:qm:volatility-models:6}
Ses trois régresseurs calculent la moyenne de la \omterm{def:qm:volatility-models:rv}{variance réalisée} sur un jour, une semaine et un mois, si bien que la prévision est une somme de composantes à mémoire courte, moyenne et longue~; cette
cascade imite l'autocorrélation à décroissance lente, à \omterm{def:qm:linear-time-series:longmem}{mémoire longue}, de la volatilité, tout en restant une régression linéaire facile à ajuster et difficile à surajuster.
\iqlookfor{Les horizons multiples comme approximation de la \omterm{def:qm:linear-time-series:longmem}{mémoire longue}.}
\end{solution}
